\section*{Chapitre \ref{ch:g4:plants-without-seeds} --- Les plantes sans graines}

\begin{solution}{exo:g4:plants-without-seeds:1}
Les \omterm{prop:g4:plants-without-seeds:damp}{mousses} et les \omterm{prop:g4:plants-without-seeds:damp}{fougères}. Elles se reproduisent par \omterm{def:g4:plants-without-seeds:spore}{spores} ---
des grains microscopiques libérés en très grand nombre.
\end{solution}

\begin{solution}{exo:g4:plants-without-seeds:2}
Une \omterm{def:g4:plants-without-seeds:spore}{spore} est microscopique et nue --- pas de réserve, pas de petite
\omterm{def:g1:plants-around-us:plant}{plante} toute prête~; une \omterm{def:g2:from-seed-to-plant:seed}{graine} porte les deux, sous une enveloppe
protectrice. Et les \omterm{def:g4:plants-without-seeds:spore}{spores} se comptent par millions, les \omterm{prop:g3:flowers-fruits-seeds:transform}{graines} par
douzaines.
\end{solution}

\begin{solution}{exo:g4:plants-without-seeds:3}
Sur le revers des frondes~: de nettes rangées de points bruns,
chacun un paquet de milliers de \omterm{def:g4:plants-without-seeds:spore}{spores} qui s'ouvre à maturité et les
laisse s'envoler comme de la poussière.
\end{solution}

\begin{solution}{exo:g4:plants-without-seeds:4}
La mousse se reproduit par \omterm{def:g4:plants-without-seeds:spore}{spores}, et les \omterm{def:g4:plants-without-seeds:spore}{spores} comme leurs jeunes
stades délicats ont besoin d'humidité --- ils sèchent et meurent
dans l'air sec. Le mur nord ombragé reste humide~; le mur sud
ensoleillé cuit.
\end{solution}

\begin{solution}{exo:g4:plants-without-seeds:5}
Faire une nouvelle \omterm{def:g1:plants-around-us:plant}{plante} directement à partir d'un morceau de la
\omterm{def:g1:plants-around-us:plant}{plante} mère --- stolon, \omterm{ex:g4:plants-without-seeds:bulbtuber}{bulbe}, \omterm{ex:g4:plants-without-seeds:bulbtuber}{tubercule}, bouture. Un seul parent.
\end{solution}

\begin{solution}{exo:g4:plants-without-seeds:6}
La \omterm{def:g1:plants-around-us:plant}{plante} envoie une longue \omterm{def:g1:plants-around-us:parts}{tige} fine --- le stolon --- sur le sol~;
là où elle touche terre, elle s'enracine et pousse des \omterm{def:g1:plants-around-us:parts}{feuilles},
devenant une petite \omterm{def:g1:plants-around-us:plant}{plante} complète~; le stolon se dessèche ensuite
et l'enfant se tient seul, identique à son parent.
\end{solution}

\begin{solution}{exo:g4:plants-without-seeds:7}
Une pomme de terre plantée (un \omterm{ex:g4:plants-without-seeds:bulbtuber}{tubercule}) germe à coup sûr par ses
«~\omterm{def:g1:the-five-senses:sense}{yeux}~», dont chacun peut donner une \omterm{def:g1:plants-around-us:plant}{plante} entière, identique à la
bonne variété plantée --- plus vite et plus sûrement que des \omterm{prop:g3:flowers-fruits-seeds:transform}{graines}.
\end{solution}

\begin{solution}{exo:g4:plants-without-seeds:8}
Réponse libre. Des \omterm{def:g1:plants-around-us:parts}{racines} blanches apparaissent sur la \omterm{def:g1:plants-around-us:parts}{tige}
immergée, en général au bout d'une à deux semaines~; la bouture est
alors prête à être plantée --- une nouvelle menthe complète, issue
d'un seul parent.
\end{solution}

\begin{solution}{exo:g4:plants-without-seeds:9}
La \omterm{def:g4:plants-without-seeds:spore}{spore} est nue et bon marché~: sans réserve ni plantule, presque
toutes se perdent, si bien qu'il faut en faire des millions. La
\omterm{def:g2:from-seed-to-plant:seed}{graine} est un coûteux nécessaire de survie, si bien qu'une douzaine
de \omterm{prop:g3:flowers-fruits-seeds:transform}{graines} bien équipées suffit. Elles font écho à la morue (des
millions, sans soins) et à l'éléphante (peu, très bien pourvus).
\end{solution}

\begin{solution}{exo:g4:plants-without-seeds:10}
Les fraisiers sont venus par stolon --- \omterm{def:g4:plants-without-seeds:vegetative}{reproduction végétative}, un
seul parent, des copies exactes. Les haricots sont venus de \omterm{prop:g3:flowers-fruits-seeds:transform}{graines}
faites par des \omterm{def:g3:flowers-fruits-seeds:parts}{fleurs}, qui mêlent les apports de deux parents (le
\omterm{def:g3:flowers-fruits-seeds:parts}{pollen} d'une \omterm{def:g3:flowers-fruits-seeds:parts}{fleur}, les \omterm{def:g4:animal-reproduction:sexual}{cellules œufs} d'une autre), si bien que
chaque pied de haricot varie un peu.
\end{solution}

\begin{solution}{exo:g4:plants-without-seeds:11}
\omterm{def:g4:plants-without-seeds:vegetative}{Reproduction végétative} --- la \omterm{def:g1:plants-around-us:tree}{branche} est une bouture naturelle. La
crue a fait le travail de sécateur de la jardinière~: elle a coupé
la pousse, l'a transportée et l'a plantée dans la vase humide, le
lit d'enracinement idéal du saule. Les rivières se bordent
elles-mêmes de saules en les cassant et en les plantant du même
geste.
\end{solution}
